三个子问题已通过同一个耦合磁模型连接起来。对于 P1，比例阻尼的 Birkhoff 表示、相容端点投影和生成函数导出维度一致的加权映射。普通中点仍是作用于完整向量场的另一种实现。非交换矩阵消元保留质量、刚度和阻尼因子的运算次序；有效模态质量由明确的归一化推导得到，不能与组装后的梁质量矩阵混用。

对于 P2，积分观测在同一坐标系内辨识保守残差和阻尼。补偿保留标称余弦耦合，且每一组控制消融均经过相同的物理力--力矩及磁场--梯度接口。实际实现的闭环方程同时保留势能误差与阻尼误差。受限的无噪声实验检验了这一完整联系；较小的拟合残差并不构成一般 CHNN 训练性能或已标定物理器件的结论。

对于 P3，在隐式中点方程内加入反馈后，在局部光滑性、稳定性和定义域假设下，可得到对同一实际闭环的收敛。误差界包含累积增量残差，因此代数求解器容差和积分步长不能互相替代为精度指标。固定辨识误差则产生相对于标称目标的另一项偏差。精确辛性要求消除非保守残差；针对比例残余阻尼的加权格式保持相应的共形结构。离散 KAM 结论进一步要求解析性、扭转、非共振和小量假设；有限时间收敛与几何诊断本身不能自动建立这些假设。

修订后的材料与磁体记录为上述结论提供一致的物理基础。端部惯量、模态尺度、输入换算和能量净空界贯穿全文。扩展至更多保留模态时，必须保持已训练控制器，同时处理完整端角和明确给出的完整模型残差。在更强反馈假设下的精确全模态实现是另一项理论构造。当前两模态轨迹不能替代逐级验证，也不能替代实际执行的能量约束器。

尚需完成的实证工作包括试样辨识、联合磁场--梯度及感知标定，以及独立全阶模型或 COMSOL 检验。有限振幅 KAM 证书则要求在声明的定义域上验证解析余项和定理特定常数。因此，本文给出方向一的格式、控制器、受控收敛分析及可复现验证序列，同时明确保留这些后续认证要求。
